\section{When the Subgraph Is Not There}
\label{sec:ch08-when-the-subgraph-is-not-there}

\index{router!what a dead subgraph looks like|(}%
Stop the Sessions service, leave the router running, and send the router a
query it can no longer answer:

\begin{minted}{text}
POST /graphql HTTP/1.1
Host: localhost:3002
Content-Type: application/json

{"query":"{ sessions { nodes { id title } } }"}
\end{minted}

\begin{minted}[fontsize=\footnotesize,breakanywhere]{text}
HTTP/1.1 200 OK
Cache-Control: no-store, no-cache, must-revalidate
Content-Type: application/json; charset=utf-8
Content-Length: 93
Vary: Accept-Encoding

{"errors":[{"message":"Failed to fetch from Subgraph 'sessions'."}],"data":{"sessions":null}}
\end{minted}

\index{router!answers 200 for a service that is not running}%
Two hundred. The service behind this graph is not running, nothing on the
machine can answer the query, and the response is a successful one. If anything
you own watches status codes to decide whether the graph is healthy, it saw a
success. This is not the router being careless: GraphQL puts errors in the body
and a transport-level code would be the wrong place for it. It is still the
thing that surprises people, and better met on a laptop than in production.

\index{router!names the subgraph and nothing else}%
Now read what it did say. \enquote{Failed to fetch from Subgraph 'sessions'.} That
is the whole error. Which of your services failed, and not one word about why:
no url, no status code, no connection error, no timeout. With one subgraph the
sentence is nearly redundant, because there is only one thing it could have
been. With three it is the first useful thing you get, and it is still only the
first, which is the problem chapter~\ref{ch:blame-routing} is about.

\index{null!the failure stopped at the nullable field}%
And look at where the null landed. \code{data.sessions} is null, not
\code{data} itself. The failure went as far up as the first nullable thing and
stopped there, which is exactly the mechanism
chapter~\ref{ch:schema-design} built the nullability rules around.
Chapter~\ref{ch:null-blast-radius} is what the same failure looks like when the
field it hits is not nullable, and the short version is that this response is
the good case.

\subsection{The Router Is Fine, Thank You}

\index{health check!passes while the graph cannot answer}%
While all that was true, both health endpoints answered:

\begin{center}
\begin{tabular}{lll}
\toprule
Request & Status & Body \\
\midrule
\code{GET /health} & 200 & \code{OK} \\
\code{GET /health/ready} & 200 & \code{OK} \\
\bottomrule
\end{tabular}
\end{center}

\index{health check!is about the process, not the graph}%
Ready, by the router's account, with a graph that cannot answer a query. That is
not a bug and I would not want it changed: the router is up, it holds a valid
execution config, and it will serve the moment something answers on 5001. The
readiness of a process and the health of a graph are different questions and
this endpoint answers the first one.

\index{health check!what a real one would have to do}%
The alternative costs more than it looks, and
\enquote{the health check should check the subgraphs} sounds obviously right for
about ten seconds. A router that reported unready whenever any subgraph was down
would take itself out of a load balancer during a deployment of any one of your
services, which is the moment you most want the rest of the graph served. The
partial-failure behavior above is the better answer, and the price of it is
that health is not the signal you want to alert on.

\index{observability!what this chapter hands chapter 20}%
So the question chapter~\ref{ch:operations} inherits is not
\enquote{why does the health check lie}. It is what you watch instead, given
that the router will report 200 at the transport, OK at the health endpoint, and
a named subgraph in an error body, all at once, for a service that is simply not
running.

\subsection{What Is Left}

\index{Part II!what the last chapter of it does}%
A binary, a config file, and a composed graph, and the endpoint a client talks
to is not one of your services any more. The subgraph is still on 5001 with its
entity route open, the router is on 3002 with two fewer fields and a plan for
every query, and the plan has one fetch in it because there is one place to
fetch from.

\index{Part II!the graph the rest of the book uses}%
Chapter~\ref{ch:second-third-subgraph} moves the speaker rows to a service that
owns them and adds the third subgraph, which is the last composition error
chapter~\ref{ch:composition} left going away for the right reason. It is also
the chapter where the plan grows a second step, and everything this one showed
you how to read starts being worth reading.
\index{router!what a dead subgraph looks like|)}%
